\subsection{Counting}

\subsubsection{The counting principle}

We often construct an outcome through a sequence of stages, making a choice at each stage.  
The \textit{counting} principle says: if there are \(r\) stages and at stage \(i\) we always have \(n_i\) options (independent of previous choices), then the total number of possible outcomes is
\[
n_1 n_2 \cdots n_r. 
\]

\bigskip\textsc{Example: outfits}

Suppose you have 4 shirts, 3 ties, and 2 jackets.  
Think of getting dressed as a 3-stage process:
\[
\text{shirt}: 4 \text{ choices}, \quad
\text{tie}: 3 \text{ choices}, \quad
\text{jacket}: 2 \text{ choices}.
\]
By the counting principle, the number of different outfits is
\[
4 \cdot 3 \cdot 2 = 24.
\]

\bigskip\textsc{Example: license plates}

Consider license plates of the form “2 letters followed by 3 digits.”  
We construct the plate sequentially.

\textit{With repetition allowed.}  
At each stage:
\[
\underbrace{26}_{\text{1st letter}}
\cdot
\underbrace{26}_{\text{2nd letter}}
\cdot
\underbrace{10}_{\text{1st digit}}
\cdot
\underbrace{10}_{\text{2nd digit}}
\cdot
\underbrace{10}_{\text{3rd digit}}.
\]
So the number of such plates is
\[
26^2 \cdot 10^3.
\]

\textit{Without repetition.}  
Now require that no letter and no digit can be used more than once.  
Then:
\[
\underbrace{26}_{\text{1st letter}}
\cdot
\underbrace{25}_{\text{2nd letter}}
\cdot
\underbrace{10}_{\text{1st digit}}
\cdot
\underbrace{9}_{\text{2nd digit}}
\cdot
\underbrace{8}_{\text{3rd digit}}.
\]
Thus, the number of such plates is
\[
26 \cdot 25 \cdot 10 \cdot 9 \cdot 8.
\]

\bigskip\textsc{Example: permutations and factorial}

Let a set have \(n\) distinct elements; we want to order them (form a permutation).  
Think of \(n\) slots and fill them one at a time:
\[
\underbrace{n}_{\text{1st slot}}
\cdot
\underbrace{(n-1)}_{\text{2nd slot}}
\cdot
\underbrace{(n-2)}_{\text{3rd slot}}
\cdots
\underbrace{1}_{\text{last slot}}.
\]
The total number of permutations is
\[
n! = n (n-1) (n-2) \cdots 2 \cdot 1.
\]

\bigskip\textsc{Example: number of subsets}

Let a set have \(n\) elements; we want to count its subsets.  
Consider the elements one by one. For each element, we have 2 options: include it in the subset or not.  
By the counting principle, the total number of subsets is
\[
2 \cdot 2 \cdot \dots \cdot 2 = 2^n,
\]
with \(n\) factors.

\textit{Sanity check for \(n = 1\).}  
If the set is \(\{a\}\), its subsets are
\[
\{a\}, \quad \varnothing,
\]
so there are 2 subsets, in agreement with \(2^1 = 2\).  
Note that when we count subsets we include both the whole set and the empty set.

\newpage
\subsubsection{A die-roll example}

We roll an ordinary six-sided die six times, and we are interested in the probability that all six results are different; denote this event by \(A\).  
We assume that all \(6^6\) possible sequences of six rolls are equally likely, so we are in a discrete uniform model and probabilities can be computed by counting.

\bigskip\textsc{Sample space}

A typical outcome is a sequence such as \((2,3,4,3,6,2)\).  
At each of the 6 rolls we have 6 possible results, so by the counting principle the total number of outcomes is
\[
|\Omega| = 6^6.
\]
Under the equally likely assumption, each sequence has probability \(1/6^6\).  
This is consistent with viewing the experiment as 6 independent rolls of a fair die, each outcome having probability \((1/6)^6 = 1/6^6\).

\bigskip\textsc{Counting the favorable outcomes}

Event \(A\) consists of all sequences in which no number is repeated, i.e., each of the numbers \(1,\dots,6\) appears exactly once.  
Such sequences are precisely the permutations of \(\{1,2,3,4,5,6\}\), and there are
\[
|A| = 6!
\]
of them.

\bigskip\textsc{Probability of all distinct rolls}

By the discrete uniform model,
\[
P(A) = \frac{|A|}{|\Omega|} = \frac{6!}{6^6}.
\]

\subsubsection{Combinations}

We start with a set of \(n\) elements and a non-negative integer \(k\).  
A \textit{combination} is a subset of the original set that has exactly \(k\) elements.  
The number of such subsets is denoted by the binomial coefficient
\[
\binom{n}{k},
\]
read as “\(n\)-choose-\(k\)”.

\bigskip\textsc{Two ways to count ordered \(k\)-tuples}

Consider a related problem: construct an ordered list of \(k\) \textit{distinct} elements chosen from the \(n\)-element set.

\textit{First method (direct counting).}  
Think of \(k\) slots and fill them one by one with distinct elements:
\[
n \cdot (n-1) \cdot (n-2) \cdots (n-k+1).
\]
This product can be written as
\[
n (n-1) \cdots (n-k+1) = \frac{n!}{(n-k)!}.
\]

\textit{Second method (choose then order).}  
First choose which \(k\) elements will appear (ignore order), then order them.
\[
\underbrace{\binom{n}{k}}_{\text{choose the set of }k\text{ elements}}
\cdot
\underbrace{k!}_{\text{permute these }k\text{ elements}}.
\]

Both methods count the same set of ordered lists, so
\[
\binom{n}{k} \, k! = \frac{n!}{(n-k)!}.
\]
Solving for \(\binom{n}{k}\) gives the standard formula
\[
\binom{n}{k} = \frac{n!}{k! \, (n-k)!},
\]
valid for integers \(n \ge 0\) and \(0 \le k \le n\).

\newpage
\bigskip\textsc{Extreme cases and the convention \(0! = 1\)}

Consider \(\binom{n}{n}\).  
From the combinatorial definition, choosing \(n\) elements from an \(n\)-element set leaves only one possibility: take all elements, so \(\binom{n}{n} = 1\).  
From the formula,
\[
\binom{n}{n} = \frac{n!}{n! \, 0!},
\]
so this equals \(1\) if we adopt the convention \(0! = 1\).

Now consider \(\binom{n}{0}\).  
Combinatorially, there is exactly one subset with zero elements: the empty set, so \(\binom{n}{0} = 1\).  
From the formula,
\[
\binom{n}{0} = \frac{n!}{0! \, n!} = 1,
\]
again consistent with \(0! = 1\).

\bigskip\textsc{A useful binomial identity}

The sum
\[
\binom{n}{0} + \binom{n}{1} + \cdots + \binom{n}{n}
\]
can be given a combinatorial interpretation.  
Here, \(\binom{n}{k}\) counts the subsets of size \(k\) of an \(n\)-element set.  
Summing over all \(k\) from \(0\) to \(n\) counts all subsets of the set.  
But the total number of subsets of an \(n\)-element set is \(2^n\).  
Therefore,
\[
\sum_{k=0}^{n} \binom{n}{k} = 2^n.
\]


\subsubsection{Binomial probabilities}

The binomial coefficients \(\binom{n}{k}\) are closely related to probabilities in repeated coin-tossing experiments.  
Consider a coin tossed \(n\) times independently.  
Let \(p\) be the probability of heads on a single toss, so the probability of tails is \(1-p\).  
We want the probability of observing exactly \(k\) heads in these \(n\) tosses.

\bigskip\textsc{Probability of a particular sequence}

Take a concrete example with \(n = 6\), and consider a specific sequence of heads and tails, e.g.\ \(\text{H T T H H H}\).  
By independence, the probability of this sequence is
\[
p \cdot (1-p) \cdot (1-p) \cdot p \cdot p \cdot p = p^4 (1-p)^2.
\]
In general, for any \textit{given} sequence with \(k\) heads and \(n-k\) tails, the probability of that sequence is
\[
p^k (1-p)^{n-k},
\]
because each head contributes a factor \(p\) and each tail a factor \(1-p\).

\bigskip\textsc{Number of \(k\)-head sequences}

We now want the probability of the event “exactly \(k\) heads in \(n\) tosses”.  
This event consists of all sequences of length \(n\) that contain exactly \(k\) heads.  
Every such sequence has the same probability \(p^k (1-p)^{n-k}\).  
Therefore, the desired probability equals
\[
\bigl(\text{number of sequences with exactly }k\text{ heads}\bigr) \times p^k (1-p)^{n-k}.
\]

To count the sequences with exactly \(k\) heads, think of the \(n\) tosses as \(n\) time slots.  
To specify a \(k\)-head sequence, we must choose which \(k\) of the \(n\) slots contain heads; the remaining \(n-k\) slots will contain tails.  
Thus, the number of such sequences is
\[
\binom{n}{k}.
\]

\bigskip\textsc{The binomial formula}

Putting everything together, the probability of exactly \(k\) heads in \(n\) independent tosses of a coin with head probability \(p\) is
\[
P(\text{exactly }k\text{ heads}) = \binom{n}{k} \, p^k (1-p)^{n-k},
\]
which is called the \textit{binomial} probability formula.



\newpage
\subsubsection{Example: Each person gets an ace}

We have a standard 52-card deck, dealt fairly to four players, 13 cards each (as in bridge).  
Outcomes are all partitions of the 52 cards into four hands of size 13, and all such partitions are assumed equally likely.  
We want the probability that each player gets \textit{exactly one ace}.

\bigskip\textsc{Counting via multinomial coefficients}

The sample space size is the number of ways to partition 52 distinct cards into hands of sizes
\(13,13,13,13\), namely the multinomial coefficient
\[
|\Omega|
= \binom{52}{13,13,13,13}
= \frac{52!}{13!\,13!\,13!\,13!}.
\]

Now count the favorable outcomes where each person gets exactly one ace.

First distribute the four aces:  
there are
\[
4 \cdot 3 \cdot 2 \cdot 1 = 4!
\]
ways to assign the four distinct aces to the four distinct players (each player gets one).  

Then distribute the remaining 48 non-aces arbitrarily, giving 12 additional cards to each player.  
The number of such partitions is
\[
\binom{48}{12,12,12,12}
= \frac{48!}{12!\,12!\,12!\,12!}.
\]

By the counting principle, the number of favorable outcomes is
\[
|A|
= 4! \cdot \binom{48}{12,12,12,12}
= 4! \cdot \frac{48!}{12!\,12!\,12!\,12!}.
\]

Thus, the desired probability is
\[
P(A)
= \frac{|A|}{|\Omega|}
= \frac{4! \cdot \dfrac{48!}{12!\,12!\,12!\,12!}}
       {\dfrac{52!}{13!\,13!\,13!\,13!}}.
\]

\bigskip\textsc{Alternative “aces first” viewpoint}

There is a faster way to obtain the same probability.  
Imagine the 52 slots corresponding to the four players’ 13-card hands.  
Place the four aces \textit{first}, one at a time, uniformly at random among the 52 slots.

The first ace can go anywhere.  

For the second ace, there are 51 remaining slots, of which 39 belong to players who do not yet have an ace (3 players \(\times\) 13 slots each).  
So the conditional probability that the second ace goes to a different player than the first is
\[
\frac{39}{51}.
\]

Assuming this has happened, two players now have aces.  
For the third ace there are 50 remaining slots; the “good” slots are those belonging to the two players without aces, i.e.\ \(2 \times 13 = 26\) slots.  
So the corresponding probability is
\[
\frac{26}{50}.
\]

Assuming all three aces have gone to different players, three players now have aces.  
For the fourth ace, there are 49 remaining slots, of which 13 belong to the last player without an ace, giving probability
\[
\frac{13}{49}.
\]

Multiplying these conditional probabilities, the overall probability that the four aces land in slots belonging to four different players is
\[
P(A)
= \frac{39}{51} \cdot \frac{26}{50} \cdot \frac{13}{49}.
\]

One can check algebraically that this expression is equal to the earlier ratio of multinomial coefficients.
